\section*{Chapitre \ref{ch:iv2:valuation-and-financial-statements} --- Valorisation et états financiers}
\markboth{\protect\omTocMark{}{Réponses~: Valorisation et états financiers}}{\protect\omTocMark{}{Réponses~: Valorisation et états financiers}}

\begin{solution}{iq:iv2:valuation-and-financial-statements:1}
Le résultat net baisse de 7.5. Le flux de trésorerie opérationnel et la trésorerie augmentent de 2.5. Les immobilisations diminuent de 10, et les bénéfices non distribués baissent de 7.5.

L'amortissement est une charge, donc le résultat avant impôt baisse de 10, et l'impôt baisse de 2.5, ce qui laisse le résultat net en baisse de 7.5. Dans le tableau des flux de trésorerie, on part du résultat net et l'on réintègre la dotation supplémentaire aux amortissements, car elle n'est pas un flux de trésorerie, ce qui donne $-7.5 + 10 = +2.5$~: le seul effet réel sur la trésorerie est l'impôt économisé. Au bilan, les actifs varient de $+2.5 - 10 = -7.5$, et les capitaux propres varient de $-7.5$ par le biais des bénéfices non distribués. Le bilan est donc équilibré.

\iqwrong{«~L'amortissement n'est pas un flux de trésorerie, donc la trésorerie ne change pas.~» La charge elle-même n'est pas un flux de trésorerie, mais elle réduit l'impôt à payer, et cela, lui, en est un.}
\iqfollow{Et si ces intérêts supplémentaires s'élevaient à 10, payés en trésorerie~?}{Le résultat net baisse toujours de 7.5, et il n'y a rien à réintégrer, donc la trésorerie baisse aussi de 7.5. Sous les US GAAP, les intérêts payés figurent dans le flux de trésorerie opérationnel, qui baisse donc lui aussi de 7.5. Les normes IFRS ont permis à une entreprise de les présenter dans le flux de financement, si bien que le flux opérationnel ne baisse pas nécessairement, et la norme IFRS 18, applicable à partir de 2027, fait du financement la règle pour la plupart des entreprises. Les passifs ne changent pas, et les capitaux propres baissent de 7.5.}
\end{solution}

\begin{solution}{iq:iv2:valuation-and-financial-statements:2}
Le résultat net est de 30, et le flux de trésorerie opérationnel est de $-10$. Le bénéfice est positif alors que la trésorerie baisse.

Le résultat avant impôt vaut $100 - 60 = 40$, l'impôt s'élève à 10, et le résultat net à 30. Pour le flux de trésorerie, on part de 30. Les créances clients augmentent de 100, montant qu'on retranche. Les stocks augmentent puis diminuent de 60, si bien que le solde est nul. Les dettes fournisseurs augmentent de 60, montant qu'on ajoute. Soit $30 - 100 + 60 = -10$, ce qui correspond à l'impôt payé. Les ventes et les achats n'ont entraîné aucun mouvement de trésorerie.

Les actifs valent 100 de créances et une trésorerie de $-10$, soit 90. Les passifs valent 60 de dettes fournisseurs, et les capitaux propres correspondent aux bénéfices non distribués de 30. Et $90 = 60 + 30$.

\iqwrong{Confondre bénéfice et trésorerie. Une entreprise qui croît à crédit peut afficher des bénéfices en hausse et se retrouver à court de liquidités, et c'est le cas de beaucoup.}
\end{solution}

\begin{solution}{iq:iv2:valuation-and-financial-statements:3}
La valeur d'entreprise vaut $1000 + 300 - 100 + 50 + 30 = 1\,280$, et le multiple est de 8.0 fois. Aucune des trois opérations ne modifie la valeur d'entreprise.

Valeur d'entreprise\index{valeur d'entreprise} est la valeur des activités pour l'ensemble de ceux qui ont un droit sur elles. Elle ajoute aux capitaux propres les autres créances sur l'entreprise et retranche la trésorerie, dont l'exploitation n'a pas besoin. Elle se rattache à l'EBITDA, grandeur calculée avant toute rémunération des apporteurs de capitaux.

Si elle emprunte 100, la dette et la trésorerie augmentent toutes deux de 100. Si elle verse un dividende de 50, la capitalisation boursière et la trésorerie baissent toutes deux de 50. Si elle rachète 100 d'actions, la capitalisation boursière et la trésorerie baissent de 100. Dans chaque cas, les activités n'ont pas changé, et c'est précisément l'intérêt de la mesure. Dans les deux derniers cas, la capitalisation boursière baisse parce que de la valeur a quitté l'entreprise. Elle ne baisse pas à cause de mauvaises nouvelles.

\iqwrong{Mélanger les détenteurs de droits~: comparer la valeur d'entreprise au résultat net, ou un cours d'action à l'EBITDA. Le numérateur et le dénominateur d'un multiple doivent concerner les mêmes détenteurs de droits.}
\end{solution}

\begin{solution}{iq:iv2:valuation-and-financial-statements:4}
Le flux de trésorerie disponible pour la firme vaut 100. Le flux de trésorerie disponible pour les capitaux propres vaut 95. On actualise le premier au coût moyen pondéré du capital, car il revient à l'ensemble des apporteurs de capitaux.

Pour la firme~: on impose le résultat opérationnel comme s'il n'y avait pas de dette, $200 \times 0.75 = 150$, on réintègre les 40 de charges non monétaires, et l'on retranche les 70 d'investissements et les 20 de besoin en fonds de roulement. On obtient 100.

Pour les capitaux propres~: on retranche les intérêts nets de leur bouclier fiscal, soit $20 \times 0.75 = 15$, et l'on ajoute les 10 d'emprunts nouveaux, ce qui donne 95. Vérification par le haut~: le résultat net vaut $(200 - 20) \times 0.75 = 135$, et $135 + 40 - 70 - 20 + 10 = 95$.

\iqwrong{Retrancher les intérêts puis actualiser au coût moyen pondéré du capital. Le taux intègre déjà le coût de la dette et son bouclier fiscal, de sorte qu'on compterait deux fois les intérêts.}
\iqfollow{Et qu'en est-il du flux de trésorerie disponible pour les capitaux propres~?}{On l'actualise au coût des capitaux propres, ce qui donne directement la valeur des capitaux propres. Avec des hypothèses cohérentes, les deux voies concordent.}
\end{solution}

\begin{solution}{iq:iv2:valuation-and-financial-statements:5}
Elle vaut $100/(0.09 - 0.03) \approx 1\,667$. Avec une croissance de 4~\%, ou avec un coût du capital de 8~\%, elle vaut 2\,000, soit 20~\% de plus. Si 100 est le flux de cette année, celui de l'année suivante est de 103, et la valeur vaut environ 1\,717.

C'est la formule de croissance de Gordon\index{modèle de croissance de Gordon}~: la valeur est égale au flux de l'année suivante divisé par l'écart entre le taux d'actualisation et la croissance. L'écart vaut ici 6 points, si bien que faire varier celui-ci d'un point représente un changement important. Une croissance de 2~\% donne $100/0.07 \approx 1\,429$, soit une baisse de 14~\%. Cette variation est moindre que la hausse de 20~\% obtenue avec une croissance de 4~\%, car la valeur est convexe par rapport à l'écart.

\iqwrong{Prendre le flux de cette année au numérateur. La formule actualise un flux qui arrive dans un an, et prendre celui de cette année sous-estime la valeur d'une fraction égale au taux de croissance.}
\iqfollow{Quel multiple cela implique-t-il~?}{Environ 16.7 fois le flux de trésorerie disponible de l'année suivante. Tout multiple cache un écart entre un taux d'actualisation et un taux de croissance, et il vaut la peine de se demander quelles valeurs ils prennent.}
\end{solution}

\begin{solution}{iq:iv2:valuation-and-financial-statements:6}
Il y a 118 actions diluées, et la valeur des capitaux propres dilués est de 2\,950. Comme l'obligation convertible est traitée comme des actions, on retire ses 200 de la dette dans le bridge menant à la valeur d'entreprise.

Pour les options, selon la méthode du rachat d'actions\index{méthode du rachat d'actions}~: si elles sont exercées, elles apportent $10 \times 10 = 100$, ce qui permet de racheter 4 actions à 25, soit un ajout net de 6. Les actions gratuites apportent 2 en totalité, puisqu'il n'y a rien à payer. L'obligation convertible est dans la monnaie, car 25 est supérieur à 20, si bien qu'elle devient $200/20 = 10$ actions. Donc $100 + 6 + 2 + 10 = 118$, et à 25 cela fait 2\,950, contre 2\,500 hors dilution.

\iqwrong{Compter les 200 deux fois, comme 10 actions dans la valeur des capitaux propres et de nouveau comme dette. Il faut n'en retenir qu'un, selon que l'obligation convertible est dans la monnaie.}
\iqfollow{Et à un cours de 15~?}{L'obligation convertible reste une dette de 200, et les options n'ajoutent que $10 - 100/15 \approx 3.33$ actions. Le nombre dépend du cours, et, dans une valorisation très serrée, cela rend le calcul circulaire.}
\end{solution}

\begin{solution}{iq:iv2:valuation-and-financial-statements:7}
Le résultat opérationnel augmente de 40~\%, passant de 20 à 28. Le résultat avant impôt augmente de 80~\%, passant de 10 à 18. Le résultat opérationnel s'annule pour un chiffre d'affaires de 75, et le résultat avant impôt pour un chiffre d'affaires de 87.5.

Chaque unité supplémentaire de chiffre d'affaires ajoute 0.8 de résultat, si bien que 10 de chiffre d'affaires supplémentaires en ajoutent 8. La base de 20 est faible, car les charges fixes absorbent l'essentiel de la marge sur coûts variables. L'effet de levier opérationnel\index{effet de levier opérationnel} est le rapport entre cette marge et le résultat opérationnel, soit $80/20 = 4$. L'effet de levier financier est de $20/10 = 2$. Ensemble, ils se multiplient pour donner 8, si bien qu'une hausse de 10~\% du chiffre d'affaires entraîne une hausse de 80~\%.

Au seuil de rentabilité, $0.8R = 60$ donne $R = 75$, et $0.8R = 70$ donne $R = 87.5$. Il suffit que le chiffre d'affaires baisse de 12.5~\% pour effacer le résultat avant impôt.

\iqwrong{Supposer que l'effet de levier ne joue que dans un sens. Le multiplicateur agit de la même façon à la baisse. Une baisse de 10~\% du chiffre d'affaires ramène le résultat avant impôt de 10 à 2.}
\end{solution}

\begin{solution}{iq:iv2:valuation-and-financial-statements:8}
Le coût des capitaux propres est de 10~\%, et le coût moyen pondéré du capital est de 8.75~\%. Avec une dette égale aux capitaux propres, le bêta des capitaux propres vaut environ 1.77, le coût des capitaux propres environ 12.84~\%, et le coût moyen pondéré du capital environ 8.67~\%.

Par le modèle d'évaluation des actifs financiers\index{modèle d'évaluation des actifs financiers}, le coût des capitaux propres vaut $4\% + 1.2 \times 5\% = 10\%$. Après impôt, la dette coûte $5\% \times 0.75 = 3.75\%$. Les pondérations sont de 80~\% et 20~\%, soit $8\% + 0.75\% = 8.75\%$.

Pour retirer l'effet de levier, $\beta_u = \beta / \left(1 + (1 - t)D/E\right) = 1.2/(1 + 0.75 \times 0.25) \approx 1.01$. Pour le rétablir à $D/E = 1$, on multiplie par $1 + 0.75 = 1.75$, ce qui donne environ 1.77. Le coût des capitaux propres vaut $4~\% + 1.77 \times 5~\%$, soit environ 12.84~\%. Le nouveau coût moyen pondéré du capital vaut $0.5 \times 12.84~\% + 0.5 \times 6~\% \times 0.75 \approx 8.67~\%$.

\iqwrong{Garder le coût des capitaux propres à 10~\% après l'ajout de dette. Cela ferait baisser le coût du capital à 7.25~\%, un déjeuner gratuit. La dette rend les capitaux propres plus risqués, ce qui restitue l'essentiel du gain.}
\end{solution}

\begin{solution}{iq:iv2:valuation-and-financial-statements:9}
L'entreprise vaut 1\,125, ses capitaux propres 625, son coût des capitaux propres 13~\%, et son coût moyen pondéré du capital environ 8.89~\%. Il n'avoisine pas 7~\%, car les capitaux propres deviennent plus chers à mesure que la dette augmente.

Sans dette, l'entreprise vaut $100/10\% = 1\,000$. Les intérêts procurent une économie d'impôt de $25\% \times 5\% \times 500$ chaque année, à perpétuité. Actualisée au taux propre de la dette, soit 5~\%, cette économie vaut $tD = 125$. C'est le bouclier fiscal\index{bouclier fiscal}, et la valeur est $1\,000 + 125$, comme l'ont montré Modigliani et Miller en présence d'impôts\index{théorème de Modigliani-Miller}. Les capitaux propres valent $1\,125 - 500 = 625$.

Les actionnaires perçoivent $100 - 500 \times 5~\% \times 0.75 = 81.25$, si bien que leur rendement vaut $81.25/625 = 13~\%$. Puis $\tfrac{625}{1125} \times 13~\% + \tfrac{500}{1125} \times 3.75~\% \approx 8.89~\%$. Et $100/8.89~\% = 1\,125$, ce qui boucle le raisonnement.

\iqwrong{Garder le coût des capitaux propres à 10~\%, ce qui donne environ 7.22~\%. L'entreprise vaudrait alors environ 1\,385. Le bouclier fiscal ne vaut que 125, et la dette ne réduit pas le risque des activités.}
\iqfollow{Doit-elle donc emprunter sans limite~?}{Pas en réalité. L'énoncé écarte le coût de détresse financière, et en pratique ce coût, ainsi que la hausse du coût de la dette, fixe la limite.}
\end{solution}

\begin{solution}{iq:iv2:valuation-and-financial-statements:10}
Elle vaut environ 2\,101, et la valeur terminale représente environ 78~\% de ce total. Avec une croissance terminale de 4~\%, elle vaut environ 2\,446. Si les flux de trésorerie arrivent en milieu d'année, la valeur est multipliée par $1.09^{1/2} \approx 1.044$, pour atteindre environ 2\,193.
\begin{center}\small
\begin{tabular}{@{}lrrrrr@{}}
Année & 1 & 2 & 3 & 4 & 5 \\ \hline
Flux de trésorerie disponible & 100.0 & 110.0 & 121.0 & 133.1 & 146.4 \\
Valeur actualisée à 9~\% & 91.7 & 92.6 & 93.4 & 94.3 & 95.2
\end{tabular}
\end{center}
Les cinq années donnent environ 467. La valeur terminale à l'année 5 vaut $146.41 \times 1.03/(0.09 - 0.03) \approx 2\,513$. Actualisée par $1.09^5 \approx 1.539$, elle vaut environ 1\,634.

Ces valeurs actualisées sont presque égales, car une croissance de 10~\% compense presque l'actualisation à 9~\%. Une somme rapide donne donc à peu près $5 \times 93$.

\iqwrong{Affirmer qu'une part terminale de 78~\% est un défaut de la méthode. C'est ce à quoi ressemblent les entreprises à longue durée de vie. Cela signifie que l'attention doit porter sur les hypothèses qui sous-tendent ce chiffre, à savoir la marge à long terme et le rendement du réinvestissement, et non sur la troisième décimale de l'année 2.}
\iqfollow{Un point de croissance terminale de plus fait passer la valeur terminale d'environ 2\,513 à environ 3\,045. Est-ce plausible~?}{Cela représente 21~\% de plus, pour un seul point de croissance. On la rapproche du réinvestissement que cette croissance exige, faute de quoi on obtient une croissance gratuite.}
\end{solution}

\begin{solution}{iq:iv2:valuation-and-financial-statements:11}
Un multiple de 10 implique une croissance d'environ 3.32~\%. Une croissance de 2~\% implique environ 8.0 fois. Un multiple de 18 exige environ 5.8~\%, ce qui est trop élevé pour durer à l'infini.

La valeur terminale vaut $\mathrm{FCF}(1 + g)/(r - g)$. En posant qu'elle est égale à $TV$ et en résolvant, on obtient
\[
g = \frac{TV \cdot r - \mathrm{FCF}}{TV + \mathrm{FCF}} = \frac{1000 \times 0.09 - 55}{1000 + 55} \approx 3.32~\%.
\]
Pour $g = 2~\%$, la valeur terminale vaut $55 \times 1.02/0.07 \approx 801$, soit environ 8.0 fois l'EBITDA. Pour 1\,800, la même formule donne $(162 - 55)/1855 \approx 5.8~\%$.

Une entreprise qui croît de 5.8~\% à perpétuité finirait par dépasser n'importe quelle économie. C'est donc ce qu'il faut dire~: un multiple de 18 suppose un taux supérieur à la croissance nominale de long terme, à moins que les marges ou les besoins en capitaux ne changent aussi.

\iqwrong{Traiter le multiple et le taux de croissance comme deux contrôles indépendants. Ce sont une seule et même hypothèse écrite deux fois~: il faut toujours convertir l'un en l'autre, et regarder ce que cela oblige à croire.}
\iqfollow{Les 55 sont-ils un bon point de départ~?}{Demander quelle part en est déprimée par les investissements de croissance de l'année 5. Une entreprise qui ralentit a besoin de moins de réinvestissement, son flux de base augmente, et la croissance implicite baisse.}
\end{solution}

\begin{solution}{iq:iv2:valuation-and-financial-statements:12}
Elle vaut 0 à un rendement de 5~\%, 1\,000 à 10~\%, et 1\,500 à 20~\%. Sans croissance, elle vaut 1\,000. La croissance n'ajoute de la valeur que lorsque le rendement du nouvel investissement dépasse le coût du capital.

Pour croître au taux $g$, il faut réinvestir une fraction $g/\mathrm{ROIC}$ du bénéfice, car le nouveau bénéfice est égal à l'investissement multiplié par le rendement. Le flux de trésorerie disponible vaut donc $\mathrm{NOPAT}(1 - g/\mathrm{ROIC})$, et
\[
V = \frac{\mathrm{NOPAT}\left(1 - g/\mathrm{ROIC}\right)}{r - g}.
\]
Avec un rendement de 20~\%, on réinvestit 25 sur 100, ce qui ajoute $25 \times 20~\% = 5$ de bénéfice, soit une croissance de 5~\%. Le flux de trésorerie vaut 75, et $75/5~\% = 1\,500$. Avec un rendement de 10~\%, le flux vaut 50, et $50/5~\% = 1\,000$, soit la même valeur que sans croissance. Croître au coût du capital ne crée donc aucune valeur. Avec un rendement de 5~\%, tout le bénéfice est réinvesti, et aucune trésorerie ne sort jamais.

\iqwrong{Utiliser un taux de croissance plus élevé pour justifier un prix plus haut, sans se demander ce qu'il faut pour l'obtenir. Un DCF avec croissance et sans réinvestissement suppose un repas gratuit.}
\iqfollow{Comment cela se relie-t-il au multiple cours/bénéfice~?}{Avec le rendement des fonds propres à la place du ROIC, le multiple prospectif vaut $(1 - g/\mathrm{ROE})/(r - g)$. Pour un rendement de 15~\%, un coût des fonds propres de 10~\% et une croissance de 5~\%, cela donne environ 13.3.}
\iqwhy{C'est la formule la plus utile pour distinguer une bonne entreprise d'une simple grande entreprise. Elle dit quand la croissance vaut d'être payée.}
\end{solution}

\begin{solution}{iq:iv2:valuation-and-financial-statements:13}
La VE sur l'EBITDA vaut 10 fois pour les deux. Le multiple cours/bénéfice vaut 16 pour $A$ et environ 11.4 pour $B$. Il diffère parce que les deux ont des niveaux d'endettement différents. Avec une dette à 12~\%, le multiple de $B$ vaut environ 28.6.

L'EBITDA vaut 120 pour les deux, et $1200/120 = 10$. L'entreprise $A$ a un résultat net de $100 \times 0.75 = 75$ et des fonds propres de 1\,200, soit 16 fois. L'entreprise $B$ verse des intérêts de 30, son résultat net vaut donc 52.5, et ses fonds propres valent 600, soit environ 11.4 fois.

La règle compare le rendement bénéficiaire sans levier au coût de la dette après impôt. Le rendement bénéficiaire de $A$ vaut $75/1200 = 6.25~\%$, et la dette coûte 3.75~\% après impôt. Remplacer des fonds propres chers par une dette moins chère augmente le bénéfice par unité de fonds propres, donc le multiple baisse. À 12~\%, la dette coûte 9~\% après impôt, ce qui est supérieur à 6.25~\%. Le résultat net vaut $(100 - 72) \times 0.75 = 21$, et le multiple augmente.

\iqwrong{Dire que $B$ est bon marché parce que son multiple cours/bénéfice est plus faible. Les deux ont exactement les mêmes activités et la même valeur d'entreprise. Le multiple plus faible de $B$ reflète simplement des fonds propres plus risqués.}
\end{solution}

\begin{solution}{iq:iv2:valuation-and-financial-statements:14}
Le flux de trésorerie disponible vaut 100 dans les deux cas. La comptabilisation en charges donne un EBITDA, un résultat opérationnel et un flux de trésorerie d'exploitation de 100, et un flux d'investissement de 0. L'immobilisation donne un EBITDA de 200, un résultat opérationnel de 180, un flux d'exploitation de 200 et un flux d'investissement de $-100$. En régime permanent, le résultat opérationnel redevient égal. Avec des dépenses qui croissent de 20~\% par an, il reste supérieur, d'environ 70 à l'année 6.

L'immobilisation transfère 100 d'une charge d'exploitation vers un investissement. L'EBITDA et le flux d'exploitation s'améliorent donc tous deux de la totalité des 100. Le résultat opérationnel n'en supporte qu'un cinquième, soit 20, et s'améliore de 80. La trésorerie, elle, n'a pas changé.

Avec 100 chaque année, au bout de cinq ans l'amortissement vaut $5 \times 20 = 100$, comme pour la comptabilisation en charges, bien que l'EBITDA reste supérieur de 100. Avec une croissance de 20~\%, les dépenses de l'année 6 valent environ 248.8 et l'amortissement environ 178.6, soit la moyenne des dépenses des cinq dernières années, et l'écart est d'environ 70. Une entreprise qui immobilise en croissance continue d'afficher davantage qu'une entreprise qui comptabilise en charges.

Selon les IFRS, la recherche est comptabilisée en charges, et le développement est immobilisé lorsqu'il satisfait à des critères précis, de sorte que deux entreprises peuvent réellement différer.

\iqwrong{Comparer les deux sur l'EBITDA ou le flux d'exploitation, sans rien ajuster. Comparer sur le flux de trésorerie disponible, ou ramener les deux à une même politique.}
\iqfollow{Qu'est-ce qui éveillerait vos soupçons~?}{Une hausse du montant immobilisé, en part du chiffre d'affaires ou des dépenses totales, sans nouveau produit pour l'expliquer.}
\end{solution}

\begin{solution}{iq:iv2:valuation-and-financial-statements:15}
Le cycle de conversion de la trésorerie\index{cycle de conversion de la trésorerie} vaut 75 jours, et le besoin en fonds de roulement net est d'environ 172.6. Une croissance de 20~\% absorbe environ 34.5, et la réduction de 15 jours de stocks libère environ 24.7.

Les créances valent $1000 \times 45/365 \approx 123.3$. Les stocks valent $600 \times 60/365 \approx 98.6$. Les dettes fournisseurs valent $600 \times 30/365 \approx 49.3$. Le besoin en fonds de roulement net vaut donc $123.3 + 98.6 - 49.3 = 172.6$. Le cycle vaut $45 + 60 - 30 = 75$.

Les délais restant inchangés, le besoin en fonds de roulement croît proportionnellement au chiffre d'affaires, de sorte qu'une croissance de 20~\% exige 20~\% de plus, soit 34.5 de trésorerie. Quinze jours de stocks valent $600 \times 15/365 \approx 24.7$.

\iqwrong{Affirmer que la croissance ne coûte pas de trésorerie. Ici, chaque 100 de chiffre d'affaires supplémentaire exige environ 17 de besoin en fonds de roulement. Une croissance rapide avec un long cycle est une raison fréquente pour laquelle des entreprises rentables manquent de liquidités.}
\iqfollow{Et si le cycle était négatif~?}{Les clients paient alors avant que les fournisseurs ne soient payés, et la croissance produit de la trésorerie. Les abonnements et certaines activités de distribution peuvent fonctionner ainsi, et cette trésorerie cesse lorsque la croissance s'arrête.}
\end{solution}

\begin{solution}{iq:iv2:valuation-and-financial-statements:16}
La direction affiche un rendement de 5~\%. En traitant la rémunération en actions comme une charge, il vaut environ 3.33~\%. Maintenir le nombre d'actions constant revient à en émettre environ 1.67 par an. À 20 fois, les deux valorisations valent 3\,000 et 2\,000, donc la plus élevée dépasse la plus faible de 50~\%.

La capitalisation boursière vaut $100 \times 30 = 3\,000$. Le rendement affiché vaut $150/3000$, et le rendement ajusté $100/3000$. Les 50 nouvelles actions émises à 30 représentent environ 1.67 action, soit une dilution d'environ 1.67~\% par an. Les racheter coûterait les 50 en trésorerie, ce qui ramène le flux de trésorerie disponible réel à 100.

Les IFRS exigent que la rémunération en actions soit comptabilisée en charge dans les comptes. Elle est réintégrée dans le flux de trésorerie, car aucune trésorerie ne sort. Mais les salariés ont tout de même été payés, et les actionnaires l'ont financé par la dilution.

\iqwrong{«~Elle ne donne lieu à aucune sortie de trésorerie, elle n'est donc pas un coût réel.~» On peut payer ce coût en trésorerie par un rachat, ou par la dilution en ne rachetant pas. L'ignorer dans les deux cas revient à valoriser l'entreprise comme si son personnel travaillait gratuitement.}
\iqfollow{Et si vous la déduisez tout en utilisant un nombre d'actions dilué~?}{Comptez les attributions futures une seule fois. Les attributions existantes entrent dans le nombre d'actions, et les attributions futures entrent comme une charge. Utiliser les deux pour une même attribution revient à compter deux fois.}
\end{solution}

\begin{solution}{iq:iv2:valuation-and-financial-statements:17}
En actions, le bénéfice par action vaut $120/115 \approx 1.043$, en hausse de 4.3~\%. En dette, il vaut $106.5/100 = 1.065$, en hausse de 6.5~\%. À un prix de 450 en actions, il tombe à environ 0.98. Une opération qui relève le bénéfice par action ne crée pas de valeur pour autant.

L'entreprise $A$ gagne 1 par action, à un multiple cours/bénéfice de 20. En actions, elle émet $300/20 = 15$ actions. En dette, les intérêts valent $300 \times 6~\% = 18$, soit 13.5 après impôt, de sorte que le résultat net vaut $120 - 13.5 = 106.5$.

La règle pour les actions consiste à comparer les multiples~: une cible achetée à 15 fois est relutive\index{accrétion et dilution} pour un acheteur à 20 fois. À 450, le prix vaut 22.5 fois, au-dessus de 20, et l'opération est dilutive. La règle pour la dette consiste à comparer les rendements~: le rendement bénéficiaire de $B$ à ce prix vaut $20/300 \approx 6.67\%$, au-dessus du coût de la dette après impôt de 4.5~\%.

\iqwrong{Lire l'accrétion comme une création de valeur. Le bénéfice par action augmente chaque fois que l'on achète un multiple plus faible, ce qui ne dit rien sur le fait que 300 soit un bon prix. Le multiple de l'acheteur peut baisser pour refléter un portefeuille plus lent ou plus risqué, et les actionnaires n'en tirent aucun profit.}
\iqfollow{Qu'est-ce qui détermine alors la valeur~?}{Le fait que le prix soit inférieur à ce que $B$ vaut pour vous, y compris les synergies, actualisé. Comparez le rendement des 300 au coût du capital de la cible.}
\end{solution}

\begin{solution}{iq:iv2:valuation-and-financial-statements:18}
La valeur des fonds propres à la sortie est d'environ 976, soit un multiple de 2.44 sur 400, et un rendement annuel d'environ 19.5~\%. À 9 fois, les fonds propres valent environ 849, soit un multiple de 2.12, et le rendement vaut environ 16.2~\%.

L'EBITDA croît jusqu'à $100 \times 1.05^5 \approx 127.6$. À 10 fois, la valeur d'entreprise vaut environ 1\,276. La dette est tombée à 300, de sorte que les fonds propres valent environ 976. Le rendement vaut $2.44^{1/5} - 1$. Un contrôle rapide~: à 15~\%, l'argent double en cinq ans environ, et 2.44 représente un peu plus qu'un doublement, donc le rendement est un peu supérieur à 15~\%.

Le gain total de 576 provient de deux sources~: 276 de la croissance de l'EBITDA à multiple constant, et 300 du remboursement de la dette. Aucun ne vient d'une variation du multiple. C'est la dette qui fait que la hausse de 28~\% de la valeur d'entreprise se traduit par une hausse de 144~\% des fonds propres.

\iqwrong{Supposer que le multiple augmentera à la sortie. Un modèle qui en a besoin repose sur un pari sur le marché, et non sur l'entreprise. Un modèle sérieux le maintient donc constant, voire à la baisse.}
\iqfollow{Qu'est-ce qui limite la dette~?}{La capacité à payer les intérêts et à rembourser une mauvaise année. Vérifier la couverture des intérêts sous une baisse de l'EBITDA, avant de regarder le rendement.}
\end{solution}

\begin{solution}{iq:iv2:valuation-and-financial-statements:19}
La dette locative et l'actif valent tous deux environ 421.2. L'année 1, l'amortissement vaut environ 84.2 et les intérêts environ 25.3, l'EBITDA augmente de 100, et le résultat avant impôt baisse d'environ 9.5. La valeur d'entreprise, dette locative comprise, vaut environ 1\,421, et l'EBITDA vaut 250, ce qui donne environ 5.68 fois, contre 6.67 auparavant. La trésorerie ne change pas. La charge totale sur 5 ans reste à 500.

La valeur actuelle de cinq paiements de 100 à 6~\% vaut 421.24. L'actif s'amortit uniformément sur 5 ans, soit $421.24/5 = 84.25$. Les intérêts de l'année 1 valent $6~\% \times 421.24 = 25.27$. Ensemble, ils font 109.52, contre un loyer de 100. Les intérêts diminuent à mesure que la dette baisse, pour atteindre 346.51 après l'année 1, puis zéro, de sorte que la charge est plus élevée au début et plus faible à la fin~: de 109.52 l'année 1 à 89.91 l'année 5.

L'EBITDA gagne la totalité du loyer, car le loyer est sorti des charges d'exploitation. Le résultat opérationnel ne gagne que $100 - 84.25 = 15.75$. Le paiement en trésorerie de 100 est le même, mais sur ce montant, $100 - 25.27 = 74.73$ compte désormais comme le remboursement d'une dette.

Il s'agit de l'IFRS 16\index{comptabilisation des contrats de location}. Selon le modèle américain, un tel contrat de location figure lui aussi au bilan, mais la charge reste en général une charge d'exploitation linéaire.

\iqwrong{Croire que l'entreprise paraît meilleure. L'EBITDA augmente, mais la dette aussi. Si l'on compare un multiple entre entreprises, les deux doivent être sur une même base, sinon le multiple plus faible n'est qu'une illusion.}
\end{solution}

\begin{solution}{iq:iv2:valuation-and-financial-statements:20}
Le résultat net vaut 30 chaque année. Le reste varie d'une année à l'autre~:
\begin{center}\small
\begin{tabular}{@{}lrrrrr@{}}
Année & 1 & 2 & 3 & 4 & 5 \\ \hline
Résultat net & 30.0 & 30.0 & 30.0 & 30.0 & 30.0 \\
Impôt décaissé & 2.5 & 2.5 & 15.0 & 15.0 & 15.0 \\
Solde d'impôt différé & 7.5 & 15.0 & 10.0 & 5.0 & 0.0 \\
Flux de trésorerie d'exploitation & 57.5 & 57.5 & 45.0 & 45.0 & 45.0
\end{tabular}
\end{center}
À 8~\%, payer plus tard au lieu de payer de façon régulière vaut environ 2.3.

Dans les comptes, le résultat avant impôt vaut $60 - 20 = 40$, la charge d'impôt vaut 10 et le résultat net vaut 30. Pour l'impôt des années 1 et 2, la déduction est de 50, donc le bénéfice imposable vaut 10 et l'impôt décaissé vaut 2.5. Pendant les années 3 à 5, il n'y a aucune déduction, donc l'impôt décaissé vaut 15. La charge de 10 dépasse l'impôt décaissé de 7.5 pendant les deux premières années, et cet écart constitue un passif, le solde d'impôt différé\index{impôt différé}. Il se résorbe ensuite de 5 par an. Selon la norme IAS 12, une entreprise comptabilise cet écart pour les différences temporelles entre la valeur comptable d'un actif et sa base fiscale.

L'impôt total payé vaut 50 dans les deux cas, soit $25~\% \times 200$. Le gain ne tient donc qu'au décalage dans le temps~: $10$ par an actualisés, moins le profil réel, donnent environ 2.3.

\iqwrong{Considérer l'impôt ainsi différé comme une économie définitive. Il s'inverse, sauf si l'entreprise continue d'acheter de nouveaux actifs. Une entreprise dont l'investissement croît peut différer l'impôt pendant de nombreuses années, et cela vaut davantage.}
\end{solution}

\begin{solution}{iq:iv2:valuation-and-financial-statements:21}
Le goodwill vaut 225. La réévaluation réduit le résultat net de 7.5 par an. Une dépréciation de 100 réduit le résultat net de 100, sans effet sur la trésorerie ni sur l'impôt.

Le goodwill\index{goodwill} est ce qui reste une fois les actifs ramenés à leur juste valeur. Les actifs nets de la cible à la juste valeur valent les capitaux propres comptables de 200, plus la réévaluation de 100, moins un passif d'impôt différé de $25\% \times 100 = 25$. Ils valent donc 275, et $500 - 275 = 225$. Vérification~: la variation des actifs de l'acquéreur vaut $+200 + 100 + 225 - 500 = +25$, ce qui égale le nouveau passif de 25.

Un amortissement supplémentaire de 10 par an réduit le résultat avant impôt de 10. Comme il n'est pas déductible, l'impôt décaissé ne change pas. La charge d'impôt comptabilisée baisse de 2.5 à mesure que l'impôt différé se résorbe, donc le résultat net baisse de 7.5.

La norme IFRS 3 affecte le prix aux actifs et passifs identifiables à leur juste valeur, et le reste va au goodwill. La norme IAS 36 soumet le goodwill à un test de dépréciation chaque année, et une perte sur goodwill n'est jamais reprise. Sous les normes IFRS comme sous les règles américaines applicables aux sociétés cotées, le goodwill n'est pas amorti.

\iqwrong{Oublier l'impôt différé et répondre 200. La réévaluation n'est pas déductible, donc l'impôt futur sera plus élevé que ne le montrent les comptes, et l'acquéreur a repris ce passif.}
\iqfollow{Que vous indique une dépréciation~?}{Que la transaction paraît désormais surévaluée, ou que l'activité s'est dégradée depuis la clôture de l'opération. Il faut donc y voir un indice sur la qualité de la gestion, et non une nouvelle perte.}
\end{solution}

\begin{solution}{iq:iv2:valuation-and-financial-statements:22}
Le prix implique une croissance constante d'environ 5.88~\% à perpétuité. Avec une croissance de 3~\% après 10 ans, il implique environ 14.3~\% par an pendant ces 10 ans. Avec une croissance de 3~\% dès aujourd'hui, l'entreprise vaut environ 1\,236.

Pour une croissance constante, il faut résoudre $3000 = 60(1 + g)/(0.08 - g)$, ce qui donne $g = 180/3060 = 1/17 \approx 5.88~\%$. Avec un taux de croissance de 3~\%, la valeur vaut $61.8/0.05 = 1\,236$.

La version à deux étapes n'admet pas de forme fermée utile, il faut donc la résoudre numériquement. À 14.3~\% pendant 10 ans, le flux de trésorerie est multiplié par environ 3.8, pour atteindre environ 228. La valeur terminale actualisée représente alors environ 72~\% du prix.

Il s'agit d'une DCF inversée\index{DCF inversée}. On ne demande pas ce que vaut l'entreprise. On demande ce qu'il faut croire pour payer ce prix. Le 1\,236 signifie qu'environ 59~\% du prix rémunère une croissance supérieure à 3~\%. Il faut alors comparer les 14.3~\% à ce qu'ont réalisé des entreprises comparables, et à ce que le réinvestissement devrait rapporter.

\iqwrong{Prendre 5.88~\% pour une prévision. Ce taux dépasse toute croissance de long terme plausible, donc la forme à croissance constante indique seulement que le marché anticipe une période de croissance plus rapide.}
\iqfollow{Et si vous croyez à 10~\%~?}{Alors les actions sont trop chères avec ces hypothèses, et il faut chercher ce que voit le marché. Il peut s'agir des marges, et non de la croissance.}
\end{solution}

\enlargethispage{-2\baselineskip}%
\begin{solution}{iq:iv2:valuation-and-financial-statements:23}
Le chiffre d'affaires vaut 213, et le résultat d'exploitation vaut 20.15, en hausse d'environ 11.9~\% par rapport à 18. Un point de ventes à périmètre constant déplace le résultat de 0.6, soit environ 3.0~\% de 20.15. Avec des ventes à périmètre constant stables, le résultat d'exploitation vaut 18.35, en hausse d'environ 1.9~\%.

Année de base~: chiffre d'affaires de 200, résultat des magasins de 30, frais centraux de 12, donc résultat d'exploitation de 18.

L'année suivante~: les magasins existants vendent $200 \times 1.03 = 206$, et les 6 de chiffre d'affaires supplémentaires ajoutent $6 \times 30~\% = 1.8$ au résultat des magasins. Les nouveaux magasins vendent chacun à un rythme annuel de $2 \times 0.7 = 1.4$, pendant un demi-an, donc $10 \times 1.4 \times \tfrac12 = 7$, et à 5~\% cela rapporte 0.35. Le résultat des magasins vaut $30 + 1.8 + 0.35 = 32.15$. Moins 12, cela donne 20.15.

Chaque point de ventes à périmètre constant représente 2 de chiffre d'affaires, et à un taux de transmission de 30~\% cela fait 0.6. Avec des ventes stables, les magasins existants gagnent leurs 30 habituels, les nouveaux magasins ajoutent 0.35, et le résultat d'exploitation vaut 18.35. Si les nouveaux magasins n'avaient pas ouvert non plus, on retrouverait exactement les 18 de l'année précédente, comme il se doit. Sur les 2.15 de croissance, 1.8 provient donc des ventes à périmètre constant, soit environ 84~\%, et les nouveaux magasins apportent peu la première année.

\iqwrong{Appliquer la marge moyenne de 15~\% à la variation des ventes à périmètre constant, ce qui ne donne que 0.9 et un résultat d'exploitation de 19.25. Les 15~\% sont calculés après toutes les charges, et une grande partie de ces charges ne varie pas avec les ventes. Une variation du chiffre d'affaires se répercute à un taux bien supérieur à la marge moyenne, à la hausse comme à la baisse.}
\iqfollow{Et ensuite~?}{Il faut demander ce que doivent rapporter les nouveaux magasins une fois matures pour qu'il vaille la peine de les ouvrir, et s'ils prennent des ventes aux magasins existants.}
\end{solution}

\begin{solution}{iq:iv2:valuation-and-financial-statements:24}
Le résultat net a progressé de 30~\%, mais la trésorerie a baissé et chaque poste du besoin en fonds de roulement a évolué dans le mauvais sens. Le flux de trésorerie disponible est passé de $+50$ à $-50$. Avec des dépenses de développement représentant la même part du chiffre d'affaires que l'année~1, le résultat net aurait été d'environ 95.5, soit une baisse depuis 100 au lieu d'une hausse de 30~\%.

Les contrôles, avec 365 jours~:
\begin{center}\small
\begin{tabular}{@{}lrr@{}}
 & Année 1 & Année 2 \\ \hline
Jours de créances & 43.8 & 63.9 \\
Jours de stocks & 60.8 & 86.2 \\
Résultat net moins flux de trésorerie d'exploitation & $-10$ & $+70$ \\
Flux de trésorerie disponible, après développements capitalisés & $+50$ & $-50$
\end{tabular}
\end{center}
Les accruals\index{qualité des résultats} de 70 représentent environ 6.2~\% des actifs moyens, ce qui est élevé. Si les créances avaient suivi les ventes, elles auraient atteint 144, donc 66 d'entre elles, soit 5.5~\% du chiffre d'affaires, sont excédentaires. Cela suggère un chiffre d'affaires comptabilisé trop tôt, ou des clients qui ne paient pas. Les stocks dépassent de 50 leur ancien ratio au coût des ventes, ce qui indique des ventes en recul ou des marchandises qui s'accumulent.

Les produits constatés d'avance ont baissé de 20. Les facturations, soit le chiffre d'affaires plus la variation des produits constatés d'avance, s'élèvent à 1\,180, soit moins que le chiffre d'affaires de 1\,200. Une partie du chiffre d'affaires de cette année est donc puisée dans l'ancien solde et n'est pas remplacée.

Les développements capitalisés représentent 5.8~\% du chiffre d'affaires, contre 2~\%. À 2~\%, ils auraient atteint 24, donc 46 de plus auraient été passés en charges, soit environ 34.5 après impôt. Je néglige l'amortissement, qui serait faible ici.

\textbf{Explications innocentes à vérifier.} Une grosse vente proche de la clôture expliquerait les créances, et les encaissements de janvier permettraient de trancher. Un lancement de produit pourrait justifier à la fois les stocks et les développements. Un gros client aux délais de paiement plus longs est un vrai changement, mais c'est un coût, non une croissance.

\textbf{Questions à poser à la direction.} Quelle part des ventes du T4 a été réalisée au cours du dernier mois~? Comment les conditions de paiement ont-elles évolué~? Qu'est-ce qui est immobilisé, et pourquoi maintenant~? Qu'a-t-on encaissé depuis la clôture de l'exercice~?

\iqwrong{Déclarer une fraude. Ce sont des signaux d'alerte, non une preuve. Chacun a une cause réelle possible. Ce que l'on doit à un recruteur, c'est un classement, un chiffre pour chacun et le contrôle qui permettrait de trancher.}
\iqwhy{Un vendeur à découvert et un investisseur de long terme lisent tous deux un tableau de ce type, et c'est cette compétence qui les distingue de celui qui ne lit que le titre.}
\end{solution}

\begin{solution}{iq:iv2:valuation-and-financial-statements:25}
L'ancien prix impliquait une reprise de 33~\%, et le nouveau prix implique environ 18.6~\%. Il faut environ 32.1~\% pour un rendement attendu de 15~\%. À 35~\%, le rendement attendu vaut environ 18.3~\%. Une nouvelle position dimensionnée selon la règle du 1~\% représenterait environ 4.8~\% du capital, contre environ 3~\% au départ. La position a perdu 0.9~\% du capital jusqu'ici. La recommandation est de conserver la position, ou de l'alourdir seulement si l'on peut expliquer pourquoi la révision à la baisse des prévisions ne remet pas en cause le scénario de reprise, et de la réduire sinon.

Chaque prix est une moyenne pondérée\index{probabilité implicite} par les probabilités de ses deux scénarios, donc $p = (P - \text{baissier})/(\text{haussier} - \text{baissier})$. Ainsi $(90 - 60)/90 = 1/3$, et $(63 - 50)/70 \approx 18.6\%$. Un rendement de 15~\% exige alors $(63 \times 1.15 - 50)/70 \approx 32.1\%$. À 35~\%, la valeur vaut $0.35 \times 120 + 0.65 \times 50 = 74.5$, et $74.5/63 - 1 \approx 18.3\%$.

À l'origine, la perte potentielle valait $30/90 = 33.3~\%$, donc une limite de perte de 1~\% correspondait à $1~\%/33.3~\% = 3~\%$ du capital. Elle a perdu $3~\% \times 27/90 = 0.9~\%$. Désormais, la perte potentielle vaut $13/63 \approx 20.6~\%$, donc la même règle autorise environ 4.8~\%, et le gain potentiel atteint environ 90.5~\%.

Le calcul invite à acheter davantage, mais tout dépend de l'hypothèse de 35~\%. Une première erreur a déjà été commise, et le nouveau scénario baissier n'est qu'une estimation. Avant d'alourdir, il faut se demander ce que la baisse dit de la reprise, si le scénario baissier constitue désormais réellement le plancher, et si l'on achèterait le titre aujourd'hui sans son historique. Tenir bon pour tenir bon ne constitue pas un argument.

\iqwrong{S'ancrer sur 90, le prix payé. Ce prix ne donne plus aucune information. L'autre erreur classique consiste à vendre parce que la position fait mal, sans point de vue sur les chiffres.}
\iqfollow{Qu'est-ce qui vous ferait changer d'avis~?}{Énoncez-le à l'avance~: par exemple, une marge du prochain trimestre passant sous un seuil annoncé, ou la perte de clients. Si cela se produit, on vend sans nouveau débat.}
\iqwhy{Chaque entretien sur les fondamentaux finit ici~: une position, une mauvaise nouvelle, et la question de savoir si l'on raisonne en probabilités et en tailles de position, ou si l'on défend ce que l'on détient déjà.}
\end{solution}
